% !TEX root = ../main.tex
\chapter{Discussion}
\label{sec:discussion}

This chapter interprets the results of \cref{sec:results}. \Cref{sec:discussion-mechanism} explains why the bias works harder under Muon, and \cref{sec:discussion-router} discusses which optimizer should train the router. \Cref{sec:discussion-overload} discusses the overload of the first MoE layer, and \cref{sec:discussion-limitations} states the limits of our results.

\section{What Raises the Work of the Bias Under Muon}
\label{sec:discussion-mechanism}
In RQ1, we expected Muon to lower the span. Under AdamW, the update of an expert grows with its load, and under Muon it does not. The feedback in which busy experts become busier (\cref{sec:background-feedback}) should therefore be weaker under Muon. This premise held, but the span rose, mostly in the first MoE layer. With the whole backbone on Muon, RQ3 and the two probes trace this rise to the inputs of the first router rather than to the updates of the experts.

Our results support the following explanation. With the backbone on Muon, the inputs of the first router share a direction, so their mean is long compared with a single input. By \cref{eq:background-offset}, the mean input adds to the logit of each expert an offset that every token shares. The more the rows of the router differ along the mean input, the more these offsets differ across the experts. Every token therefore prefers the same experts, and the bias has to cancel this preference, which widens its span.

Three results support this explanation. First, across the balanced runs of RQ3, the span of the first MoE layer followed the score offset much more closely than the input cosine (\cref{fig:results-geometry}). Second, the span of this layer and the shared direction both depended on the optimizer of the backbone, not on that of the router. Third, we centered the input of the first router, hiding the shared direction from it. This removed the early overload in both probes without changing the loss. Over the first 3,000 steps, centering also removed 79\% of the excess span on the small model and all of it on the large one.

The input cosine alone explained little of the span because the score offset has two factors. The first is the length of the mean input, which the backbone sets. The second is how much the rows of the router differ along it, which depends on how the router is trained. At half the rate of Muon, the first factor stayed as it was, but the second fell from 3.0 to 1.4 times its value under AdamW (\cref{sec:results-offset}). With the fast bias, which leaves the rate as it was, it fell to 1.05 times (\cref{app:results-rq3}). In both cases, the early overload was also below twice the fair share, against about four times at the tuned rate. This suggests that the rows move apart mostly while the loads are uneven. The score offset may therefore feed itself: it causes the overload, and the overload moves the rows further apart.

We see one way in which the overload could move the rows apart, although our experiments did not test it. Recall that the gradient of a row of the router comes only from the tokens routed to its expert (\cref{sec:background-feedback}). While the bias lags behind the score offset, the favored experts receive more tokens than the others. Their rows therefore receive more of the gradient along the shared direction, and they move apart from the other rows along it. A larger rate would move them further. So would a router on AdamW, since its update acts more on the shared direction (\cref{sec:results-offset}).

The shared direction does not explain the span with the routed experts alone on Muon. This partition gave the largest span in RQ3, although its input cosine was only 0.15 (\cref{sec:results-rq3-parts}). The tuning runs of the control suggest that the size of the updates plays a role. Recall that the control gives the update of every expert the same size, as Muon does, but without the orthogonalization (\cref{sec:method-optimizers}). Its span was higher than that of AdamW at every rate above its best, although with one run per rate (\cref{app:results-tuning}). Updates of the same size for every expert may therefore raise the span. This is the opposite of what we expected in RQ1, and our results do not show how it happens.

The main question that our experiments leave open is why only the whole backbone on Muon lengthens the mean input. Wang et~al.~\cite{wang2026myth} describe the shared direction in the hidden states of trained MoE models. Li \cite{li2026spectralcap} reports a related effect in the first MoE layer of a model trained with Muon, but without a run with AdamW\@. Our results add that, in the first MoE layer, the strength of this direction depends on the optimizer of the backbone. Under Muon, the direction also weakened with depth (\cref{sec:results-rq1}). This fits a note of Wang et~al.: under Muon, the similarity of the hidden states did not grow with depth.

The shared direction may also explain why our runs without the bias agree only in part with those of Wang et~al.~\cite{wang2026depth}. Averaged over the layers, Muon kept the loads more even than AdamW, as they found. In the first MoE layer, however, the loads collapsed further under Muon (\cref{sec:results-rq3-unbalanced}). In this layer, the inputs shared the direction under Muon, and without the bias nothing canceled the resulting offset.

\section{Where to Put the Router}
\label{sec:discussion-router}
The recipes of current frontier models disagree on the optimizer of the router (\cref{tab:related-router}). Only the Qwen3.8-Flash-Next report gives a reason: Muon made the router unstable early in training and gave no clear gain later. At the scale of our models, our results point the other way. With the backbone on Muon, moving the router to AdamW made three things worse. First, the warmup churn was higher on both model sizes, and at every learning rate of the stress test of RQ2. Second, the early overload of the first MoE layer was higher on the small model, at every seed. Third, the loss was higher at every seed, by 0.002 on the small model and by 0.004 on the large one.

When the inputs of the router share a direction, the gradient of every one of its rows has a large component along it. A step along this direction changes the logit of an expert by about the same amount for every token. The common share measures how much of the update of the router acts along this direction. AdamW normalizes its update entry by entry and keeps much of this component. Muon makes the 32 singular values of its update nearly equal, which holds the common share near $1 / \sqrt{32} \approx 0.18$ (\cref{sec:method-readings}). On the small model, over steps 301 to 1,262, the common share was about 0.48 with the router on AdamW against 0.19 under Muon (\cref{sec:results-offset}). Much of the update of a router on AdamW therefore raises or lowers the logit of an expert for all tokens at once, which fits its higher churn and overload. Our results do not show why the loss with the router on AdamW was higher, but the reason is not the early overload of the first MoE layer. In the probe of the large model, centering removed this overload under both Muon partitions, yet the loss with the router on AdamW stayed 0.005 to 0.007 above that of Muon (\cref{app:results-large}).

Our results and the Qwen3.8-Flash-Next report may disagree for several reasons. First, its models are much larger. The model of its stress tests, for example, has 25B parameters, about fifty times as many as our large model \cite{qiu2026qwennext}. Second, the report shows no data that compare the two optimizers on the router. Third, it does not state its balancer or report the loads of its experts. We therefore cannot tell whether the early instability that it describes is the warmup churn that we measure.

Our results also point to one possible drawback of the router on Muon. In the stress test of RQ2, its largest logit grew with the rate. At 16 times the tuned rates, it reached 98 on average over the seeds and 170 at one seed, whereas a router on AdamW stayed below 36 at every rate (\cref{app:results-stress}). In a router that normalizes its logits with a softmax, large logits cause rounding errors, which the router z-loss prevents (\cref{sec:background-instability}). Our sigmoid router does not normalize its logits. Since no run diverged and the runs had too few loss spikes, we cannot tell whether the large logits harmed training.

With the backbone on AdamW, the optimizer of the router mattered little for the first MoE layer and for the loss. Moving the router alone to Muon changed the span of this layer by $+0.004$ and the loss by $-0.005$. The optimizer of the router therefore matters mainly when the backbone is on Muon, the case on which the recipes disagree. Srivastava \cite{srivastava2026lossfree} trained the attention and the experts with Muon in every run, and the router with Adam or with Muon. The router on Muon was better conditioned (\cref{sec:related-muon-balance}), and the loss was nearly the same. This agrees with the small cost in loss of the router on AdamW.

\section{The Overload of the First MoE Layer}
\label{sec:discussion-overload}
Our models have no capacity limit (\cref{sec:background-moe}). Every token therefore reached its selected experts, even when one expert received several times its fair share. In the probes, removing the early overload left the loss as it was. In other settings, however, the overload would have costs that our runs do not show. When the experts of a layer sit on different devices, a few busy experts slow down training while the others are underused \cite{lepikhin2021gshard}. Other models give each expert a capacity, a fixed limit on the number of tokens that it processes in a batch, and drop the tokens beyond it \cite{fedus2022switch}. In such settings, an early overload of four to seven times the fair share for hundreds of steps would cost time or dropped tokens. So would the overload that remained to the end of training on the large model, two to three times the fair share. We did not measure either cost.

Three changes lowered the early overload, each with a limit. First, centering the input of the first router removed the early overload in both probes, together with most of the extra span. We tested the centering, however, only over the first 3,000 steps and in one layer. Second, a bias four times faster lowered the early overload of Muon at every seed, and over whole runs it removed about half of the excess span. In the probe of the large model at seed 0 (\cref{app:results-large}), however, the load churn was still three times that of AdamW, although lower than with the usual bias. The likely reason is that each larger step of the bias moved more tokens at once. This is the trade-off that the Kimi K3 report names: a fixed step of the bias either adapts slowly or makes the loads oscillate \cite{kimi2026k3}.

Third, the balance loss of DeepSeek-V3 lowered the score offset in every arm and the early overload of both Muon partitions. By Proposition~2 of Wang et~al.~\cite{wang2026myth}, a balance loss teaches the router to suppress the shared direction. At $\alpha = 10^{-4}$, our balance loss did this at most in part: the input cosine and the common share did not change in the same direction at every seed.

On the large model, the first MoE layer reached a state that the small model had not shown. At seed 0 of the probe, with the backbone on Muon, the scores of this layer sank toward zero (\cref{sec:results-large}). The bias therefore made most of the choices. Under both Muon partitions, more than half of the selected experts would not have been selected by the scores alone, against a quarter under AdamW (\cref{app:results-large}). Many tokens were also close to a tie between their last selected and first unselected experts. Each step of the bias therefore moved many of them at once, and the loads of this layer did not settle. The only loss spikes of the large model came with the router on AdamW, and all of them fell where the span of this layer was at or near its ceiling. The DeepSeek-V4 report also links its loss spikes to the routing (\cref{sec:background-instability}). Ten spikes in three runs, however, are too few to establish such a link.

Other reports also find the first MoE layers the hardest to balance. Sigma-MoE-Tiny found this under AdamW, and Motif 3 under Muon (\cref{sec:related-balancer}). Our results add that Muon on the backbone makes the first MoE layer harder to balance. A recipe that trains the backbone with Muon and balances the loads with the bias should therefore watch the first MoE layer most closely.

\section{Limitations}
\label{sec:discussion-limitations}
Our models have 55M and 108M active parameters, one to three orders of magnitude fewer than the models whose recipes motivate our questions. They train on 1B and 3B tokens, whereas DeepSeek-V4 was trained on more than 32T tokens \cite{deepseek2026v4}. The effects in the first MoE layer grew from the small model to the large one. The input cosine rose from 0.4 to 0.6 under Muon and 0.7 with the router on AdamW, and the early overload of Muon from 4 to 6 times the fair share. The MoE Routing Testbed also warns that routing methods that perform alike in small models often differ in large ones, although its own main finding held at 35 times the size \cite{falke2026testbed}. Our results therefore describe a mechanism and its size at our scale, not a recipe for a production model. We cannot say whether the overload of the first MoE layer would settle in a run a thousand times longer.

We train one architecture on one dataset, FineWeb-Edu. The architecture has a sigmoid router, selects 4 of 32 routed experts, adds a shared expert, places one dense block before the first MoE layer, and sets no capacity limit. A model with more dense blocks before its first MoE layer, or with none, may pass the shared direction to its routers differently. The sigmoid also bounds the scores, and this gives the span its ceiling (\cref{sec:background-span}). DeepSeek-V4, by contrast, computes the score of a logit $z$ as $\sqrt{\log(1 + e^{z})}$, which has no upper bound \cite{deepseek2026v4}. On our large model, the span of the first MoE layer reached the ceiling at some steps. With unbounded scores, it could have grown further.

Our results also rest on one balancer. We used the sign rule of \cref{eq:bias-update} at two speeds, $u$ and $4u$, and we added the balance loss at a single weight, $\alpha = 10^{-4}$. The balancers of Kimi K3 and Trinity Large (\cref{sec:related-balancer}) may respond differently to the shared direction.

The learning rates of the small model were tuned with one seed and 0.5B tokens, and then used with 1B tokens. With 1B tokens, the tuned rate of Muon was not its best, since half of it reached a lower loss at every seed. Most of the rise also came with this higher rate (\cref{sec:results-rq3-rate}). We did not check the rate of AdamW again with 1B tokens. The large model was tuned at its own size, and its rates were then scaled to 3B tokens (\cref{sec:method-tuning}). We did not check whether a lower rate of Muon would also train better there. We took the other hyperparameters from published recipes or set them by convention, for example the weight decay, the warmup, and the clipping. Of these, we checked only $\beta_2$ of AdamW, with one run at 0.98 instead of 0.95.

Each comparison rests on three seeds. If two arms do not differ, their difference has the same sign at all three seeds with probability $1/4$. We also report many comparisons without a correction for their number. Three things limit this risk. First, we fixed the comparisons and their expected signs before the runs. Second, we tested every prediction made from earlier results on new runs. Third, the main differences held at every seed in several experiments. The rise, for example, did so in RQ1, in RQ3, in the bias-speed experiment, and on the large model. The exploratory analyses and the probes, in contrast, were chosen after we had seen the runs, and the probes cover only the first 3,000 steps. The metrics of the routing are logged every 10 steps, so changes between two logs are missed. Finally, loss spikes have no standard definition, so our rule is one choice among several (\cref{sec:method-stability}).
